% !TeX root = ../main.tex
\section{Группы, порядки элементов и перестановки}

\task{55.16}{Тождество $x^2=e$}
\utv{}{Если $x^2=e$ для каждого $x\in G$, то $G$ абелева.}{
 Из $x^2=e$ следует $x^{-1}=x$. Применяя это также к $ab$, получаем
 \[ab=(ab)^{-1}=b^{-1}a^{-1}=ba.\]
 Здесь порядок множителей при взятии обратного меняется;
 коммутативность заранее не используется.
}

\task{55.21}{Изоморфизмы $\Z_4\to\Z_5^*$}
Любой гомоморфизм из $\Z_4=\langle\overline1\rangle$
определяется образом $\overline1$. Изоморфизм должен переводить
его в элемент порядка $4$. В $\Z_5^*$ такими элементами являются
только $2$ и $3$:
\[
2^2=4,\quad 2^3=3,\quad 2^4=1;
\qquad 3^2=4,\quad 3^3=2,\quad 3^4=1\pmod5.
\]
Следовательно, имеются ровно два изоморфизма:
\[\varphi_2(\overline k)=\overline{2^k},\qquad
  \varphi_3(\overline k)=\overline{3^k}.\]
Корректность обеспечивается равенствами $2^4=3^4=1$:
замена $k$ на $k+4t$ не меняет результат.
\begin{center}
\begin{tabular}{c|cccc}
$k$ & $0$ & $1$ & $2$ & $3$\\\hline
$\varphi_2(k)$ & $1$ & $2$ & $4$ & $3$\\
$\varphi_3(k)$ & $1$ & $3$ & $4$ & $2$
\end{tabular}
\end{center}

\task{55.29}{Группа аффинных преобразований}
Пусть $F$ --- поле и
\[G=\{(a,b):a\in F^*,\ b\in F\},\qquad
 (a,b)(c,d)=(ac,ad+b).\]
\utv{}{Это группа, изоморфная группе функций $x\mapsto ax+b$,
 $a\ne0$, относительно композиции.}{
 Замкнутость следует из $ac\ne0$. Ассоциативность проверяется непосредственно:
 \begin{align*}
 ((a,b)(c,d))(u,v)&=(acu,acv+ad+b),\\
 (a,b)((c,d)(u,v))&=(acu,a(cv+d)+b).
 \end{align*}
 Нейтральный элемент --- $(1,0)$, обратный ---
 \[(a,b)^{-1}=(a^{-1},-a^{-1}b).\]
 Умножение в обоих порядках даёт $(1,0)$.
 Зададим $\varphi(a,b)(x)=ax+b$. Тогда
 \[\varphi(a,b)\bigl(\varphi(c,d)(x)\bigr)=acx+ad+b
 =\varphi((a,b)(c,d))(x).\]
 Сюръективность следует из определения группы функций.
 Для инъективности достаточно подставить $x=0$ и $x=1$:
 равенство функций даёт сначала равенство $b$, затем равенство $a$.
}
\editnote{Функции $x\mapsto ax+b$ называются аффинными;
 при $b\ne0$ они не являются линейными отображениями пространства $F$.}

\task{56.8}{Произведения коммутирующих элементов}
Пусть $xy=yx$, $\ord(x)=m$, $\ord(y)=n$.
\utv{Пункт а}{Если $\gcd(m,n)=1$, то $\ord(xy)=mn$.}{
 Имеем $(xy)^{mn}=x^{mn}y^{mn}=e$, поэтому порядок произведения
 делит $mn$. Пусть $(xy)^t=e$. Тогда $x^t=y^{-t}$.
 Возведение этого равенства в степень $m$ даёт $y^{tm}=e$,
 откуда $n\mid tm$, а значит $n\mid t$. Аналогично $m\mid t$.
 Из взаимной простоты следует $mn\mid t$. Минимальный такой $t$
 равен $mn$.
}
\utv{Пункт б}{Существуют целые $k,l$ такие, что
 $\ord(x^ky^l)=\lcm(m,n)$.}{
 Для каждого простого $p$ сравним показатели $v_p(m)$ и $v_p(n)$.
 Положим
 \[
 a=\prod_{v_p(m)\ge v_p(n),\ v_p(m)>0}p^{v_p(m)},
 \qquad
 b=\prod_{v_p(n)>v_p(m)}p^{v_p(n)}.
 \]
 Каждый простой множитель попадает ровно в одно произведение;
 поэтому $a\mid m$, $b\mid n$, $\gcd(a,b)=1$ и
 $ab=\lcm(m,n)$. По формуле для порядка степени
 \[\ord(x^{m/a})=a,\qquad \ord(y^{n/b})=b.\]
 Эти степени коммутируют. Пункт а даёт требуемое для
 $k=m/a$, $l=n/b$. Пустое произведение понимается как $1$.
}
\editnote{В пункте б исходник необоснованно возвращался к случаю
 взаимно простых $m,n$. Здесь доказан общий случай.}

\task{56.18}{Период конечной группы}
\Def{Период, или экспонента, конечной группы $G$ --- наименьшее
 положительное число $d(G)$ такое, что $g^{d(G)}=e$ для всех $g\in G$.}
\utv{Пункт а}{
 $d(G)=\lcm\{\ord(g):g\in G\}$ и $d(G)\mid |G|$.
}{
 Условие $g^s=e$ для всех $g$ равносильно тому, что $s$
 делится на порядок каждого элемента. Наименьшее положительное
 такое число есть указанный НОК. По теореме Лагранжа все эти
 порядки делят $|G|$, следовательно, их НОК также делит $|G|$.
}
\utv{Пункт б}{Если $G$ абелева, в ней есть элемент порядка $d(G)$.}{
 Перечислим элементы $G$ как $g_1,\dots,g_N$. Положим $h_1=g_1$.
 Если построен $h_i$ порядка $\lcm(\ord(g_1),\dots,\ord(g_i))$,
 применим пункт б задачи 56.8 к $h_i$ и $g_{i+1}$.
 Они коммутируют, поэтому существуют степени, произведение которых
 имеет порядок $\lcm(\ord(h_i),\ord(g_{i+1}))$.
 Возьмём это произведение за $h_{i+1}$. В итоге
 $\ord(h_N)=d(G)$.
}

\task{3.2}{Разложение в независимые циклы}
Для перестановки
\[\sigma=\begin{pmatrix}1&2&3&4&5&6&7\\3&7&6&5&1&2&4\end{pmatrix}\]
последовательно прослеживаем образы:
\[1\mapsto3\mapsto6\mapsto2\mapsto7\mapsto4\mapsto5\mapsto1.\]
Поэтому $\sigma=(1\,3\,6\,2\,7\,4\,5)$.
Для второго примера $i\mapsto n+i\mapsto i$ при $1\le i\le n$, так что
\[
\begin{pmatrix}1&\cdots&n&n+1&\cdots&2n\\n+1&\cdots&2n&1&\cdots&n\end{pmatrix}
=(1\,n+1)(2\,n+2)\cdots(n\,2n).
\]
\editnote{Первый разобранный в PDF пункт подписан «в»;
 сохранена именно приведённая там перестановка.}

\task{3.4, б}{Произведение перестановок}
Пусть $\alpha=(13)(57)(246)$, $\beta=(135)(24)(67)$.
При композиции справа налево получаем
\[
\begin{array}{c|rrrrrrr}
i&1&2&3&4&5&6&7\\\hline
\beta(i)&3&4&5&2&1&7&6\\
\alpha(\beta(i))&1&6&7&4&3&5&2
\end{array}
\]
Следовательно, $\alpha\beta=(2\,6\,5\,3\,7)$; точки $1,4$ неподвижны.

\task{3.7}{Чётность перестановок}
Цикл длины $k$ раскладывается в $k-1$ транспозиций:
\[(i_1\,i_2\,\dots\,i_k)
 =(i_1\,i_k)(i_1\,i_{k-1})\cdots(i_1\,i_2).
\]
Его знак равен $(-1)^{k-1}$: цикл чётен при нечётном $k$.
Произведение $k$ независимых транспозиций
$(i_1i_2)\cdots(i_{2k-1}i_{2k})$ имеет знак $(-1)^k$.

\task{3.11}{Инверсии и соседние транспозиции}
\Def{Инверсия в строке $a_1,\dots,a_n$ --- пара $i<j$ с $a_i>a_j$.}
\utv{}{Если перестановка имеет $k$ инверсий, то минимальное число
 соседних транспозиций $(q\,q+1)$ в её разложении равно $k$.}{
 Правое умножение на $(q\,q+1)$ меняет местами соседние значения
 $a_q,a_{q+1}$. Сравнения с любым третьим значением в сумме
 остаются прежними; меняется только статус пары $(q,q+1)$.
 Поэтому число инверсий меняется ровно на $1$.

 Если строка не возрастает, в ней есть соседнее убывание:
 иначе $a_1<a_2<\cdots<a_n$, то есть строка уже равна $1,\dots,n$.
 Меняем местами соседние значения $a_q>a_{q+1}$; число инверсий
 уменьшается на $1$. Через ровно $k$ шагов получим тождественную
 перестановку. Обращая последовательность шагов, получаем
 разложение исходной перестановки в $k$ соседних транспозиций.

 За $r$ таких шагов число инверсий может измениться не более чем на $r$.
 Чтобы перейти от $0$ к $k$, нужно $r\ge k$.
}
\editnote{При $a_q>a_{q+1}$ обмен убирает инверсию;
 в исходнике это направление изменения было перепутано.}

\task{3.16}{Образующие знакопеременной группы}
\utv{Пункт а}{Любая чётная перестановка --- произведение тройных циклов.}{
 Разложим её в чётное число транспозиций и сгруппируем их по парам.
 Для каждой пары возможны три случая. Одинаковые транспозиции
 сокращаются. При одной общей точке
 \[(xz)(xy)=(xyz).\]
 Для четырёх различных точек
 \[(xy)(ab)=(xay)(xab).\]
 Последнее равенство проверяется по образам $x,y,a,b$:
 правая часть меняет $x$ с $y$ и $a$ с $b$.
 Так каждая пара заменяется нулём, одним или двумя тройными циклами.
}
\utv{Пункт б}{При $n\ge3$ группа $A_n$ порождается циклами
 $(123),(124),\dots,(12n)$.}{
 Обозначим $s_i=(12i)$, $3\le i\le n$, и
 $K=\langle s_3,\dots,s_n\rangle$.
 Все образующие чётны, так что $K\le A_n$.
 При различных $i,j\ge3$ сопряжение даёт
 \[s_i s_j s_i^{-1}=(2ij)\in K.\]
 Сопрягая $(2ij)$ тем же $s_i$, получаем $(i1j)$;
 его обратный --- $(1ij)$. Наконец, при попарно различных
 $a,b,c\ge3$ имеем
 \[s_a(2bc)s_a^{-1}=(abc).\]
 Итак, получены тройные циклы всех типов: с $1,2$;
 с одной из этих точек; без обеих точек.
 Обратные циклы также лежат в $K$.
 По пункту а это все образующие $A_n$, значит $K=A_n$.
 Для $n=3$ достаточно единственного цикла $(123)$.
}

\task{3.23}{Перестановка как произведение двух инволюций}
\utv{}{Для любой $\sigma\in S_n$ существуют $\alpha,\beta\in S_n$
 с $\alpha^2=\beta^2=e$ и $\sigma=\alpha\beta$.}{
 Сначала рассмотрим один цикл длины $l$. Занумеруем его точки
 остатками $j\in\Z_l$ так, что цикл действует как $j\mapsto j+1$.
 Положим $\alpha_l(j)=-j$, $\beta_l(j)=-j-1$.
 Тогда $\alpha_l^2(j)=\beta_l^2(j)=j$, а
 \[\alpha_l\beta_l(j)=-(-j-1)=j+1.\]
 Разложим $\sigma$ в независимые циклы и на каждом зададим
 такие две перестановки. Произведения всех $\alpha_l$ и всех
 $\beta_l$ имеют квадрат $e$, поскольку их носители не пересекаются.
 Их произведение действует на каждом цикле как $\sigma$.
}
\editnote{Транспозиции в разложении
 $(a_1a_l)\cdots(a_1a_2)$ не независимы: у всех общая точка $a_1$.
 Поэтому предложенное в исходнике разделение на «чётные» и «нечётные»
 множители не доказывало утверждение.}

\task{56.3}{Порядки конкретных элементов}
\textbf{б)} Перестановка
$\left(\begin{smallmatrix}1&2&3&4&5&6\\2&3&4&5&1&6\end{smallmatrix}\right)
=(12345)$ имеет порядок $5$.

\textbf{г)} Число $z=(1-i)/\sqrt2=e^{-\pi i/4}$ удовлетворяет
$z^t=1$ ровно при $-\pi t/4\in2\pi\Z$, то есть при $8\mid t$.
Поэтому $\ord(z)=8$.

В PDF также приведена строка ответов ко всем пунктам задачи.
Ниже она сохранена и пояснена; условия остальных пунктов восстановлены
по приложенному задачнику.
\begin{center}
\begin{tabular}{c|ccccccccc}
Пункт & а & б & в & г & д & е & ж & з & и\\\hline
Порядок & $6$ & $5$ & $12$ & $8$ & $4$ & $8$ & $2$ & $3$ & $n$
\end{tabular}
\end{center}
\textbf{а)} Перестановка $(123)(45)$ имеет порядок $\lcm(3,2)=6$.
\textbf{в)} $-\sqrt3/2+i/2=e^{5\pi i/6}$ имеет порядок $12$,
поскольку $5t$ должно делиться на $12$.
\textbf{д)} Матрица циклического сдвига четырёх координат имеет порядок $4$.
\textbf{е)} Для
$A=\left(\begin{smallmatrix}0&i\\1&0\end{smallmatrix}\right)$
имеем $A^2=iI$, $A^4=-I$, $A^8=I$; нечётная степень не скалярна,
а чётные степени равны $i^rI$, поэтому меньшего порядка нет.
\textbf{ж)} Для
$A=\left(\begin{smallmatrix}-1&a\\0&1\end{smallmatrix}\right)$
имеем $A^2=I$, $A\ne I$.
\textbf{з)} Для
$A=\left(\begin{smallmatrix}0&-1\\1&-1\end{smallmatrix}\right)$
выполняется $A^2+A+I=0$, откуда $A^3=I$; $A\ne I$.
\textbf{и)} Верхнетреугольная матрица с попарно различными диагональными
элементами $\lambda_1,\dots,\lambda_n$, составляющими все корни
$n$-й степени из $1$, диагонализуема: её характеристический многочлен
имеет простые корни, а собственные векторы для различных собственных
значений линейно независимы. Она подобна $\diag(\lambda_1,\dots,\lambda_n)$.
Её $n$-я степень равна $I$, и среди собственных значений есть примитивный
$n$-й корень, поэтому меньшая положительная степень не равна $I$.

\task{56.6}{Число элементов порядка $6$}
\textbf{б)} Рассмотрим обратимые диагональные матрицы
$\diag(a,b)$ над $\C$. Порядок матрицы равен НОК порядков
$a,b$. Если он делит $6$, оба числа --- корни шестой степени из $1$.
Пусть $N(d)$ --- число матриц, удовлетворяющих $A^d=I$.
Тогда $N(d)=d^2$, поскольку каждую диагональную координату
можно выбрать $d$ способами.
Среди $36$ матриц с $A^6=I$ исключаем порядки $1,2,3$:
\[36-4-9+1=24.\]
Прибавление $1$ возвращает матрицу $I$, вычтенную дважды.

\textbf{г)} В $A_5$ нет элементов порядка $6$.
В $S_5$ такой порядок возможен только при типе циклов $(3)(2)$,
а его знак равен $(-1)^2(-1)=-1$.
В PDF также записаны ответы: в $\C^*$ таких элементов $2$
(примитивные шестые корни), в $S_5$ ---
$\binom53\cdot2=20$ (выбор трёх точек и одного из двух тройных циклов).

\task{56.10}{Порядок перестановки}
\utv{Пункт а}{Порядок нечётной перестановки чётен.}{
 Если $\ord(\sigma)=t$, то
 $1=\sgn(e)=\sgn(\sigma^t)=(-1)^t$, откуда $2\mid t$.
}
\utv{Пункт б}{Если длины независимых циклов равны
 $l_1,\dots,l_r$, то $\ord(\sigma)=\lcm(l_1,\dots,l_r)$.}{
 На носителе $i$-го цикла степень $\sigma^t$ действует как
 степень именно этого цикла; другие циклы там тождественны.
 Поэтому $\sigma^t=e$ равносильно одновременному выполнению
 $l_i\mid t$ для всех $i$. Минимальный положительный $t$ --- их НОК.
}

\task{Т.1}{Когда перестановка является квадратом?}
\thm{Критерий квадратного корня}{Перестановка является квадратом
 другой перестановки тогда и только тогда, когда для каждой
 чётной длины число её циклов этой длины чётно.}{
 Квадрат цикла действует прыжком на две позиции.
 Если длина $l$ нечётна, получится один цикл той же длины,
 поскольку $\gcd(2,l)=1$. Если $l=2m$, то
 \[(a_1\dots a_{2m})^2
 =(a_1a_3\dots a_{2m-1})(a_2a_4\dots a_{2m}).\]
 Поэтому цикл чётной длины в квадрате возникает только в паре
 с циклом той же длины. Это доказывает необходимость.

 Для достаточности строим корень. У цикла $c$ нечётной длины $l$
 корнем служит $c^{(l+1)/2}$, так как
 $(c^{(l+1)/2})^2=c^{l+1}=c$.
 Каждую пару циклов одинаковой чётной длины
 $(a_1\dots a_l)$ и $(b_1\dots b_l)$ объединяем в цикл
 \[d=(a_1b_1a_2b_2\dots a_lb_l).
 \qquad d^2=(a_1\dots a_l)(b_1\dots b_l).\]
 Все построенные корни имеют независимые носители.
 Их произведение при возведении в квадрат даёт исходную перестановку.
}

\task{56.28}{Элемент порядка $2$ в группе чётного порядка}
\utv{}{В конечной группе чётного порядка есть элемент порядка $2$.}{
 Множество $G\setminus\{e\}$ имеет нечётное число элементов.
 Разобьём элементы с $g\ne g^{-1}$ на пары $\{g,g^{-1}\}$.
 Если бы все элементы попали в такие пары, число элементов
 было бы чётным. Поэтому есть $g\ne e$ с $g=g^{-1}$.
 Тогда $g^2=e$ и $\ord(g)=2$.
}
